% Part: proof-theory
% Chapter: normalization
% Section: permutations

\documentclass[../../../include/open-logic-section]{subfiles}

\begin{document}

\olfileid{pt}{nor}{per}

\olsection{Permutation Conversions}

Cuts in \Log{N1i} are segments that end in the major premise of
!!a{elimination} rule. One case in the normalization proof is to
reduce the length of maximal cuts, and with it the cut length of the
!!{proof}. A cut of length $>1$ can end in many ways depending on
whether the last !!{formula} occurrence of~$!A$ in the cut is the
conclusion of (\Elim{\lor} or \Elim{\lexists}) and depending on the
!!{elimination} inference it is the major premise of.

For instance, if the inferences involved are \Elim{\lor} and
\Elim{\land}, $!A \ident !B \lor !C$ and there is a
sub-!!{proof}~$\delta_1$ that ends in:
\[
\AxiomC{}
\RightLabel{$\delta_2$}
\DeduceC{$!D \lor !E$}
\AxiomC{$\Discharge{!D}{x}$}
\RightLabel{$\delta_3$}
\DeduceC{$!B \land !C$}
\AxiomC{$\Discharge{!E}{y}$}
\RightLabel{$\delta_4$}
\DeduceC{$!B \land !C$}
\DischargeRule{\Elim{\lor}}{x\,y}
\TrinaryInfC{$!B \land !C$}
\RightLabel{\Elim{\land}[1]}
\UnaryInfC{$!B$}
\DisplayProof
\]
The last !!{formula} occurrence~$!A_n$ in the cut is the conclusion
of~\Elim{\lor}, the lower $!B \land !C$. The second-to-last
!!{formula} occurrence~$!A_{n-1}$ is one of the minor premises
of~\Elim{\lor}.

To reduce the length of this cut, we can switch the order
(``permute'') of the \Elim{\lor} and \Elim{\land} inferences. That is,
we replace the subproof~$\delta_1$ in~$\delta$ by the following:
\[
\AxiomC{}
\RightLabel{$\delta_2$}
\DeduceC{$!D \lor !E$}
\AxiomC{$\Discharge{!D}{x}$}
\RightLabel{$\delta_3$}
\DeduceC{$!B \land !C$}
\RightLabel{\Elim{\land}[1]}
\UnaryInfC{$!B$}
\AxiomC{$\Discharge{!E}{y}$}
\RightLabel{$\delta_4$}
\DeduceC{$!B \land !C$}
\RightLabel{\Elim{\land}[1]}
\UnaryInfC{$!B$}
\DischargeRule{\Elim{\lor}}{x\,y}
\TrinaryInfC{$!B$}
\DisplayProof
\]
The cuts that ended in the premise of \Elim{\land} below \Elim{\lor}
now end in the premise of one of the \Elim{\land} inferences above the
minor premises of \Elim{\lor}, i.e., the length of each has decreased
by~$1$.

Moreover, any cut that lies entirely within $\delta_2$, $\delta_3$, or
$\delta_4$, or entirely outside~$\delta_1$ in~$\delta$ is unchanged:
it still involves the same !!{formula}s, goes through the same
inferences, so in particular has the same rank and length as the
corresponding cut in~$\delta$. So the cut rank of the new !!{proof} has
not increased, and the cut length of the new !!{proof} is less than that
of the old !!{proof}.

\begin{prob}
After a permutation conversion, the length of at least one cut is
reduced by~$1$. If we permute \Elim{\land} across an~\Elim{\lor},
there are actually two cut segments the length of which decrease, so
in this case the cut length of the !!{proof} decreases by at
least~$2$. Can the cut length of !!a{proof} decrease by more than~$2$
by a single such permutation? Can the cut \emph{rank} decrease?
\end{prob}

The process of moving !!a{elimination} rule above an \Elim{\lor} or
\Elim{\lexists} rule is called a \emph{permutation conversion}. The
patterns for some of the rule pairs is given in
\olref{tab:permutation-conversions}. (The remaining are left as
exercises.) 

% \begin{table}
% \begin{multline*}
% \shoveleft{\AxiomC{$!D \lor !E$}
% \AxiomC{$\Discharge{!D}{x}$}
% \shortDeduce
% \DeduceC{$!B \land !C$}
% \AxiomC{$\Discharge{!E}{y}$}
% \shortDeduce
% \DeduceC{$!B \land !C$}
% \DischargeRule{\Elim{\lor}}{x\,y}
% \TrinaryInfC{$!B \land !C$}
% \RightLabel{\Elim{\land}[1]}
% \UnaryInfC{$!B$}
% \DisplayProof}\\
% \shoveright{\longrightarrow\ 
% \AxiomC{$!D \lor !E$}
% \AxiomC{$\Discharge{!D}{x}$}
% \shortDeduce
% \DeduceC{$!B \land !C$}
% \RightLabel{\Elim{\land}[1]}
% \UnaryInfC{$!B$}
% \AxiomC{$\Discharge{!E}{y}$}
% \shortDeduce
% \DeduceC{$!B \land !C$}
% \RightLabel{\Elim{\land}[1]}
% \UnaryInfC{$!B$}
% \DischargeRule{\Elim{\lor}}{x\,y}
% \TrinaryInfC{$!B$}
% \DisplayProof}
% \\
% \shoveleft{\AxiomC{$!D \lor !E$}
% \AxiomC{$\Discharge{!D}{x}$}
% \shortDeduce
% \DeduceC{$!B \lif !C$}
% \AxiomC{$\Discharge{!E}{y}$}
% \shortDeduce
% \DeduceC{$!B \lif !C$}
% \DischargeRule{\Elim{\lor}}{x\,y}
% \TrinaryInfC{$!B \lif !C$}
% \AxiomC{$!B$}
% \RightLabel{\Elim{\lif}}
% \BinaryInfC{$!C$}
% \DisplayProof}\\
% \shoveright{\longrightarrow\ 
% \AxiomC{$!D \lor !E$}
% \AxiomC{$\Discharge{!D}{x}$}
% \shortDeduce
% \DeduceC{$!B \lif !C$}
% \AxiomC{$!B$}
% \RightLabel{\Elim{\lif}}
% \BinaryInfC{$!C$}
% \AxiomC{$\Discharge{!E}{y}$}
% \shortDeduce
% \DeduceC{$!B \lif !C$}
% \AxiomC{$!B$}
% \RightLabel{\Elim{\lif}}
% \BinaryInfC{$!C$}
% \DischargeRule{\Elim{\lor}}{x\,y}
% \TrinaryInfC{$!C$}
% \DisplayProof}
% \\
% \shoveleft{\AxiomC{$!D \lor !E$}
% \AxiomC{$\Discharge{!D}{x}$}
% \shortDeduce
% \DeduceC{$!B \lif !C$}
% \AxiomC{$\Discharge{!E}{y}$}
% \shortDeduce
% \DeduceC{$!B \lif !C$}
% \DischargeRule{\Elim{\lor}}{x\,y}
% \TrinaryInfC{$!B \lif !C$}
% \AxiomC{$!B$}
% \RightLabel{\Elim{\lif}}
% \BinaryInfC{$!C$}
% \DisplayProof}\\
% \shoveright{\longrightarrow\ 
% \AxiomC{$!D \lor !E$}
% \AxiomC{$\Discharge{!D}{x}$}
% \shortDeduce
% \DeduceC{$!B \lif !C$}
% \AxiomC{$!B$}
% \RightLabel{\Elim{\lif}}
% \BinaryInfC{$!C$}
% \AxiomC{$\Discharge{!E}{y}$}
% \shortDeduce
% \DeduceC{$!B \lif !C$}
% \AxiomC{$!B$}
% \RightLabel{\Elim{\lif}}
% \BinaryInfC{$!C$}
% \DischargeRule{\Elim{\lor}}{x\,y}
% \TrinaryInfC{$!C$}
% \DisplayProof}
% \\
% \shoveleft{\AxiomC{$!D \lor !E$}
% \AxiomC{$\Discharge{!D}{x}$}
% \shortDeduce
% \DeduceC{$\lforall[x][!B(x)]$}
% \AxiomC{$\Discharge{!E}{y}$}
% \shortDeduce
% \DeduceC{$\lforall[x][!B(x)]$}
% \DischargeRule{\Elim{\lor}}{x\,y}
% \TrinaryInfC{$\lforall[x][!B(x)]$}
% \RightLabel{\Elim{\lforall}}
% \UnaryInfC{$!B(t)$}
% \DisplayProof}\\
% \shoveright{\longrightarrow\ 
% \AxiomC{$!D \lor !E$}
% \AxiomC{$\Discharge{!D}{x}$}
% \shortDeduce
% \DeduceC{$\lforall[x][!B(x)]$}
% \RightLabel{\Elim{\lforall}}
% \UnaryInfC{$!B(t)$}
% \AxiomC{$\Discharge{!E}{y}$}
% \shortDeduce
% \DeduceC{$\lforall[x][!B(x)]$}
% \RightLabel{\Elim{\lforall}}
% \UnaryInfC{$!B(t)$}
% \DischargeRule{\Elim{\lor}}{x\,y}
% \TrinaryInfC{$!B(t)$}
% \DisplayProof}
% \end{multline*}
% \caption{Some permutation conversions for \Log{N1i}.}
% \ollabel{tab:permutation-conversions}
% \end{table}

{\small
\begin{longtable}{@{}l@{}}
\AxiomC{$!D \lor !E$}
\AxiomC{$\Discharge{!D}{x}$}
\shortDeduce
\DeduceC{$!B \land !C$}
\AxiomC{$\Discharge{!E}{y}$}
\shortDeduce
\DeduceC{$!B \land !C$}
\DischargeRule{\Elim{\lor}}{x\,y}
\TrinaryInfC{$!B \land !C$}
\RightLabel{\Elim{\land}[1]}
\UnaryInfC{$!B$}
\DisplayProof
$\longrightarrow$\\[6ex]
\quad\AxiomC{$!D \lor !E$}
\AxiomC{$\Discharge{!D}{x}$}
\shortDeduce
\DeduceC{$!B \land !C$}
\RightLabel{\Elim{\land}[1]}
\UnaryInfC{$!B$}
\AxiomC{$\Discharge{!E}{y}$}
\shortDeduce
\DeduceC{$!B \land !C$}
\RightLabel{\Elim{\land}[1]}
\UnaryInfC{$!B$}
\DischargeRule{\Elim{\lor}}{x\,y}
\TrinaryInfC{$!B$}
\DisplayProof
\\[10ex]
\AxiomC{$!D \lor !E$}
\AxiomC{$\Discharge{!D}{x}$}
\shortDeduce
\DeduceC{$!B \lif !C$}
\AxiomC{$\Discharge{!E}{y}$}
\shortDeduce
\DeduceC{$!B \lif !C$}
\DischargeRule{\Elim{\lor}}{x\,y}
\TrinaryInfC{$!B \lif !C$}
\AxiomC{$!B$}
\RightLabel{\Elim{\lif}}
\BinaryInfC{$!C$}
\DisplayProof
$\longrightarrow$\\[6ex]
\AxiomC{$!D \lor !E$}
\AxiomC{$\Discharge{!D}{x}$}
\shortDeduce
\DeduceC{$!B \lif !C$}
\AxiomC{$!B$}
\RightLabel{\Elim{\lif}}
\BinaryInfC{$!C$}
\AxiomC{$\Discharge{!E}{y}$}
\shortDeduce
\DeduceC{$!B \lif !C$}
\AxiomC{$!B$}
\RightLabel{\Elim{\lif}}
\BinaryInfC{$!C$}
\DischargeRule{\Elim{\lor}}{x\,y}
\TrinaryInfC{$!C$}
\DisplayProof
\\[10ex]
\AxiomC{$!D \lor !E$}
\AxiomC{$\Discharge{!D}{x}$}
\shortDeduce
\DeduceC{$!B \lif !C$}
\AxiomC{$\Discharge{!E}{y}$}
\shortDeduce
\DeduceC{$!B \lif !C$}
\DischargeRule{\Elim{\lor}}{x\,y}
\TrinaryInfC{$!B \lif !C$}
\AxiomC{$!B$}
\RightLabel{\Elim{\lif}}
\BinaryInfC{$!C$}
\DisplayProof
$\longrightarrow$\\[6ex]
\AxiomC{$!D \lor !E$}
\AxiomC{$\Discharge{!D}{x}$}
\shortDeduce
\DeduceC{$!B \lif !C$}
\AxiomC{$!B$}
\RightLabel{\Elim{\lif}}
\BinaryInfC{$!C$}
\AxiomC{$\Discharge{!E}{y}$}
\shortDeduce
\DeduceC{$!B \lif !C$}
\AxiomC{$!B$}
\RightLabel{\Elim{\lif}}
\BinaryInfC{$!C$}
\DischargeRule{\Elim{\lor}}{x\,y}
\TrinaryInfC{$!C$}
\DisplayProof
\\[10ex]
\AxiomC{$!D \lor !E$}
\AxiomC{$\Discharge{!D}{x}$}
\shortDeduce
\DeduceC{$\lforall[x][!B(x)]$}
\AxiomC{$\Discharge{!E}{y}$}
\shortDeduce
\DeduceC{$\lforall[x][!B(x)]$}
\DischargeRule{\Elim{\lor}}{x\,y}
\TrinaryInfC{$\lforall[x][!B(x)]$}
\RightLabel{\Elim{\lforall}}
\UnaryInfC{$!B(t)$}
\DisplayProof
\quad$\longrightarrow$\\[8ex]
\quad\AxiomC{$!D \lor !E$}
\AxiomC{$\Discharge{!D}{x}$}
\shortDeduce
\DeduceC{$\lforall[x][!B(x)]$}
\RightLabel{\Elim{\lforall}}
\UnaryInfC{$!B(t)$}
\AxiomC{$\Discharge{!E}{y}$}
\shortDeduce
\DeduceC{$\lforall[x][!B(x)]$}
\RightLabel{\Elim{\lforall}}
\UnaryInfC{$!B(t)$}
\DischargeRule{\Elim{\lor}}{x\,y}
\TrinaryInfC{$!B(t)$}
\DisplayProof
\\
\caption{Some permutation conversions for \Log{N1i}.}
\ollabel{tab:permutation-conversions}
\end{longtable}
}

\begin{prob}
Give the permutation conversions for \Elim{\lor} followed by
\Elim{\lor} or \Elim{\lnot}.
\end{prob}

\begin{prob}
Give the permutation conversions for \Elim{\lexists} followed by
\Elim{\exists}. Explain why no eigenvariable condition is violated in
the latter.
\end{prob}

We should convince ourselves that when applying a permutation
conversion, we don't perhaps violate an eigenvariable condition.
Consider
\[
\AxiomC{}
\RightLabel{$\delta_2$}
\DeduceC{$!D \lor !E$}
\AxiomC{$\Discharge{!D}{x}$}
\RightLabel{$\delta_3$}
\DeduceC{$\lexists[x][!B(x)]$}
\AxiomC{$\Discharge{!E}{y}$}
\RightLabel{$\delta_4$}
\DeduceC{$\lexists[x][!B(x)]$}
\DischargeRule{\Elim{\lor}}{x\,y}
\TrinaryInfC{$\lexists[x][!B(x)]$}
\AxiomC{$\Discharge{!B(c)}{z}$}
\RightLabel{$\delta_5$}
\DeduceC{$!C$}
\DischargeRule{\Elim{\lexists}}{z}
\BinaryInfC{$!C$}
\DisplayProof
\]
If we replace this by
\[
\AxiomC{}
\RightLabel{$\delta_2$}
\DeduceC{$!D \lor !E$}
\AxiomC{$\Discharge{!D}{x}$}
\RightLabel{$\delta_3$}
\DeduceC{$\lexists[x][!B(x)]$}
\AxiomC{$\Discharge{!B(c)}{z}$}
\RightLabel{$\delta_5$}
\DeduceC{$!C$}
\DischargeRule{\Elim{\lexists}}{z}
\BinaryInfC{$!C$}
\AxiomC{$\Discharge{!E}{y}$}
\RightLabel{$\delta_4$}
\DeduceC{$\lexists[x][!B(x)]$}
\AxiomC{$\Discharge{!B(c)}{z}$}
\RightLabel{$\delta_5$}
\DeduceC{$!C$}
\DischargeRule{\Elim{\lexists}}{z}
\BinaryInfC{$!C$}
\DischargeRule{\Elim{\lor}}{x\,y}
\TrinaryInfC{$!C$}
\DisplayProof
\]
we might worry that the eigenvariable~$c$ might occur in, say,~$!D$.
After all, the assumption $!D$ is discharged above the
original~$\Elim{\lexists}$ and so in the original !!{proof} it does not
occur in an assumption that's open above the major premise of
the~\Elim{\exists} inference, and so there the eigenvariable condition
is satisfied. But in the new !!{proof}, $!D$~is not discharged until after
the~\Elim{\lexists} inferences, and so the eigenvariable condition
would be violated. Recall, however, that we assume our !!{proof}s are
regular: every eigenvariable is the eigenvariable of a single
inference and only occurs above the minor premise of \Elim{\lexists}
in the !!{proof}. In particular, $c$~cannot occur in $\delta_3$
or~$\delta_4$.

This example shows something else: Note that the new !!{proof} has two
copies of~$\delta_5$. This means that if $\delta_5$ contains a maximal
cut, we have \emph{not} reduced the cut length! Thus we must take care
that we apply a permutation conversion only when we know that there is
no maximal cut in $\delta_5$. We will ensure this by always selecting
a \emph{topmost} maximal cut, that is a maximal cut such that the
sub-!!{proof} (in this case,~$\delta_1$) contains no maximal cut that
ends above the last inference in the sub-!!{proof}. This rules out the
possibility that there are other maximal cuts in~$\delta_5$ (or indeed
in $\delta_2$, $\delta_3$, or~$\delta_4$).

\end{document}
